%%\documentclass[t,handout]{beamer}
\documentclass{beamer}


\usepackage[utf8]{inputenc}
\usepackage[portuguese]{babel}
\usepackage[tight]{subfigure}
\usepackage{graphicx}
\usepackage{color}
\usepackage{url}
% \usepackage{listings}

%Pacote de listagem de c�digo
\usepackage{listings}
\lstset{numbers=left, stepnumber=1, firstnumber=1,
numberstyle=\tiny, extendedchars=true, breaklines=true, frame=tb,
basicstyle=\footnotesize, stringstyle=\ttfamily,
showstringspaces=false }


 \usetheme{Frankfurt} %LEGAL     !!!
% \usetheme{Madrid}     %LEGAL/L%IMPO/COM CAIXA     (sem barra de desenvolvimento)

% \usetheme{Antibes} %NAO
% \usetheme{Berlin} %PODE SER...     (BARRA DE DESENVOLVIMANTO)
% \usetheme{Berkeley}     %FEIO
% \usetheme{Boadilla} %TUDO BRANCO...
% \usetheme{Copenhagen}     %NAO
% \usetheme{Darmstadt} %LEGAL!     !!!
% \usetheme{Dresden}     %LEGAL/LIMPO/SEM CAIXA     (sem caixa fica ruim...)

% \usetheme{Goettingen}     %FEIO DEMAIS!
% \usetheme{Ilmenau} %LEGAL (forte candidato)
% \usetheme{JuanLesPins} %BACANA
% \usetheme{Luebeck}     %FEIO

% \usetheme{Malmoe}     %FEIO
% \usetheme{Warsaw} %NAO...
% \usetheme{Seattle}
% \usetheme{CambridgeUS}
% \usetheme{Singapore}

% \usecolortheme[RGB={130,35,150}]{structure}
% \usecolortheme[RGB={33,33,94}]{structure}
\usecolortheme[RGB={134,153,188}]{structure}
\setbeamertemplate{footline}[frame number]
\setbeamertemplate{navigation symbols}{}

\author[L. Medeiros]{Leonardo Medeiros}
\date{\today}
\institute[IFAL, UFCG, UFAL]{IFAL, UFCG, UFAL}
    \title{Wireless Sensor Networks for Healthcare: A Survey}
\logo{\includegraphics[width=0.2\linewidth]{img/logo.png}}
\subtitle{Hande Alemdar, Cem Ersoy}

\begin{document}

\begin{frame}
  \titlepage
\end{frame}

% \section{Roteiro}
% \AtBeginSection[]
{\frame{
\frametitle{Roteiro}
\tableofcontents
}
}

\section{Introdução}
\begin{frame}{Pervasive Healthcare}
  \begin{block}{Dados}
  O contínuo envelhecimento populacional mundial ocorrido nas últimas décadas tem se tornado um dos principais problemas dos países desenvolvidos. Previsões do Population Reference Bureau ~\cite{aging05} informam que nos próximos 20 anos, 20 $\%$ da população terá acima dos 65 anos de idade.
  \end{block}
\end{frame}

\subsection{Pervasive Healthcare}
\begin{frame}{Pervasive Healthcare}
  \begin{block}{}
   O uso da computação pervasiva aplicada ao auxílio ao tratamento de saúde é denominado de \textit{pervasive healthcare} ~\cite{ah07}. 
  \end{block}

  \begin{block}{}
  Essa tecnologia pode fornecer informação baseada no contexto e na rotina do paciente provendo um monitoramento contínuo do mesmo. 
  \end{block}
\end{frame}

\begin{frame}{Monitoramento Contínuo} 
  \begin{block}{}<1->
   O monitoramento poderá prover alertas em situações iminentes de agudização ou crise. Tratando com antecedência e trazendo menos injúria ao paciente ~\cite{alemdar}.
  \end{block}
% 
%   \begin{block}{}<2->
%    Minimiza a necessidade de acompanhantes que auxiliem: pacientes crônicos, idosos  e crianças de berço ~\cite{alemdar}.
%   \end{block}

  \begin{block}{}<2->
   O monitoramento constante irá aumentar a detecção prévia de situações de emergência para pacientes \textit{homecare} em risco ou que possuem níveis cognitivos e físicos comprometidos ~\cite{stan05}.
  \end{block}
\end{frame}


% \begin{frame}{Tratamento Domiciliar}
%  \begin{block}{}<1->
%     Dentre os benefício do uso da \textit{pervasive healthcare} tem-se a melhoria do tratamento domiciliar do paciente ~\cite{aarh10}, a realização do monitoramento remoto ~\cite{lae04}, além da possibilidade de visualizar a situação clínica do paciente através da opinião de profissionais.  
%  \end{block}
% \end{frame}


\subsection{Tecnologias}
\begin{frame}{Tecnologias Existentes}
  \begin{block}{}
  Pesquisadores da área de computação, eletrônica, medicina estão juntos pensando nas possibilidades propiciadas pela computação pervasiva aplicada a saúde. 
  \end{block}

  \begin{block}{Larga-Escala (Grandes Áreas)}
  A importância de integrar tecnologias wireless de larga escala como: 3G, Wi-Fi Mesh, WiMAX com a telemedicina têm endereçados diversas pesquisas por alguns pesquisadores.
  \end{block}  
\end{frame}

\begin{frame}{Tecnologias Existentes}
  \begin{block}{Pequena-Escala (Pequenas Áreas)}
   Algumas melhorias para esse monitoramento podem ser efetuados por tecnologias de pequenas áreas ou até mesmo da distância do corpo humano como tecnologias: RFID(Identificação de Rádio-Freqüencia), Bluetooth, Zigbee e redes sensores sem-fio.
  \end{block}

  \begin{block}{Aplicações Baseadas em Contexto}
  Através da integração das tecnologias \textit{wireless} de larga e pequena escala é possível prover aplicações sensíveis ao contexto monitorando o paciente onde quer que ele esteja ~\cite{witeleme}.
  \end{block}  
\end{frame}

\begin{frame}{Facilidade de acesso à informação}
    \begin{block}{}
     Com o avanço das tecnologias e dos sistemas de informações médicas, torna-se possível que os profissionais da saúde tenham acesso à informação do paciente, bem como submeter seu próprio parecer clínico a respeito do estado de saúde atual desses ~\cite{mug03}. 
    \end{block} 
\end{frame}

\begin{frame}{Protótipos e Produtos Comerciais}
    \begin{block}{}
      Já existem alguns protótipos e aplicações comerciais para atendimento à saúde que fazem uso de:
	\begin{itemize}
	 \item Body Area Network (BAN)
	 \item Personal Area Network (PAN)
	 \item Informações vitais do paciente em tempo-real visualizadas em Smartphones ou PDAs
	\end{itemize}
    \end{block} 
\end{frame}


\begin{frame}{}
    \begin{block}{Overview de um Wireless Sensor Network Application for Healthcare}
      \center \includegraphics[height=1.8 in]{../imagens/overwsn.png}
    \end{block}
\end{frame}


% \subsection{Futuros Cenários}
% \begin{frame}{Monitoramento Baseado no Paciente}
% Given the importance of the subject, there are already
% several applications and prototypes on the subject. For
% example, some of them are devoted to continuous moni-
% toring for cognitive disorders like Alzheimer’s, Parkinson’s
% or similar cognitive diseases. Some focus on fall detection,
% posture detection and location tracking and others make
% use of biological and environmental sensors to identify pa-
% tients’ health status. There is also significant research effort
% in developing tiny wireless sensor devices, preferably inte-
% grated into fabric or other substances and be implanted in
% human body.
% \end{frame}

\section{Benefícios}
\subsection{}
\begin{frame}{Benefícios}
\begin{block}{}
   A identificação de situações de emergência como um ataque cardíaco ou morte súbita dentro de poucos segundos ou até mesmo minutos pode ser o suficiente para salvar vidas. Por esse motivo o monitoramento de pacientes em tempo real e ações reativas providas pelos sistemas pervasivos podem ser um dos principais benefícios. 
\end{block}
\end{frame}


\begin{frame}{Avanços da Tecnologia e Redução de Custos}
O avanço da tecnologia dos equipamentos eletrônicos têm reduzido o custo e tornado possível o sensoreamento a custos baixíssimos. Conjuntamente com as tecnologias de RFID que já está madura e têm seu custo inexpressivo, é possível ``etiquetar'' usuários, dispositivos e todo o ambiente trazendo precisão as soluções pervasivas.
\end{frame}

\begin{frame}{Identificação de Contexto}
\begin{block}{Identificação de Contexto}
  A identificação de contexto é outro benefício provido pelos sistemas pervasivos aplicados a saúde. O \textit{Context-awareness} permite entender as condições ambientais que as pessoas monitoradas foram submetidas identificando o ambiente (temperatura, pressão além dos sinais vitais capturados pelos sensores). 
\end{block}

\begin{alertblock}{Informação do Contexto}
  A informação do contexto auxilia a entender e identificar padrões ainda não conhecidos, tornando mais precisas as inferências a respeito das situações monitoradas.
\end{alertblock}
\end{frame}

\begin{frame}{WSN for Healthcare}
    \begin{block}{Possíveis Usuários}
      \begin{itemize}[<+->]
       \item \textbf{Crianças:} bebês e crianças dependentes que necessitem de um monitoramento constante;
       \item \textbf{idosos e pacientes crônicos:} esse grupo incluem as pessoas que possuem dificuldades cognitivas ou alguma disfunção no coração, respiração, etc. Idosos com esses sintomas ou que sejam sucetíveis a morte súbita;
      \item \textbf{acompanhantes:} familiares, babás e acompanhantes de idosos, crianças ou pacientes crônicos;
      \item \textbf{profissionais de saúde:} Profissionais como médicos, enfermeiros e demais profissionais responsáveis pelo status clínico dos pacientes e que sejam capazes de dar respostas imediatas em casos de emergência. 
      \end{itemize}
    \end{block}
\end{frame}

\section{Design da Aplicação}
\subsection{}
\begin{frame}{} 
    \begin{block}{Considerações de Design de Aplicações Pervasivas Aplicadas à Saúde}
      \center \includegraphics[height=2.4 in]{../imagens/design.png}
    \end{block}
\end{frame}


% \subsection{Body Area Network}
\begin{frame}{Body Area Network}
    \begin{block}{}
      Body Area Network (BAN) é uma rede \textit{ad hoc} que crianças, idosos e pacientes crônicos podemo carregar em seus corpos.
    \end{block}

    \begin{block}{Exemplos:}
     Etiquetas RFID, Sensor de Eletrocardiograma (ECG), acelerômetro, medidor de temperadura são exemplos de componentes de uma BAN.
    \end{block}
\end{frame}

% \begin{frame}{Principais Problemas}
%     \begin{block}{}<1->
%      Um dos principais problemas de uma BAN é o consumo de energia, troca de baterias e recarga. Por conta disso a eficiência dos protocolos de Mediun Access Control (MAC) no que se refere a eficiência energética é crítica ~\cite{ome07}.      
%     \end{block}
% 
%     \begin{block}{}<2->
%      O estudo de Bhatia et. al. ~\cite{bha07}, demonstra que as tecnologias como Bluetooth e Wi-Fi falham no que se refere ao suporte a uma energia eficiente, desde que apenas conseguem oferecer apenas duas semanas em execução quando fazem uso de uma ``pilha de relógio''.
%     \end{block}
% \end{frame}

\begin{frame}{Principais Problemas de uma BAN}
    \begin{block}{}
      \begin{itemize}[<+->]
       \item Consumo de energia, troca de baterias e recarga. Eficiência no protocolo MAC ~\cite{ome07};
       \item tecnologias como Bluetooth e Wi-Fi falham no que se refere ao suporte a uma energia eficiente ~\cite{bha07};
       \item estudos ~\cite{ren06} investigam os efeitos biológicos causados pela transmissão de rádio freqüencia dos nós sensores e propõe um algoritmo para reduzir tais efeitos;
       \item mobilidade e portabilidade devem er considerados para o design dos sensores  ~\cite{seo07}. Sendo necessário integrar com outras tecnologias de rede como RFID  e Near Field Communication (NFC);
       \item sistemas de tempo-real disponíveis para a comunicação pode ser crítico.
      \end{itemize}
    \end{block}
\end{frame}


\begin{frame}{Problemas de Segurança de uma BAN}
    \begin{block}{}
      Problemas de segurança também devem ser considerados, pois os sinais vitais armazenados são confidenciais.
      \begin{itemize}[<+->]
       \item No trabalho de Bao ~\cite{bao05}, foi desenvolvido uma autenticação baseada nos sinais vitais. A informação extraída dos sinais fisiológicos são extraídas para gerar uma identidade para autênticação mútua;
       \item Davgtas ~\cite{dag08}, usa um algoritmo de autenticação para transmitir informações clínicas pelos sensores sobre o corpo como ECG e dispositivos móveis do paciente.
      \end{itemize}
    \end{block}
\end{frame}


% \section{Personal Area Network}

\begin{frame}{Personal Area Network}
  \begin{block}{}
      O subsistema é composto de um ambiente de sensores despostos ao redor do ambiente ou dispositovs móveis pertencentes ao paciente.  
  \end{block}

  \begin{block}{}
    O ambiente de sensores como leitores de RFID, câmeras, som, pressão, temperatura, luminosidade, humidade. Auxiliam em prover a informação contextual do paciente e do ambiente monitorado.
  \end{block}
\end{frame}

% \subsection{Aplicações}
\begin{frame}{Aplicações}
  \begin{block}{}
    \begin{itemize}[<+->]
     \item \textbf{Personal Assistant:} comunicação com três devices médicos: glicosímetro, bomba de insulina e sensor de glicose contínuo. Através de um PDA, o médico consegue controlar os dados do paciente numa aplicação sensível a contexto ~\cite{sae09};     
    \end{itemize}
  \end{block}
\end{frame}

\begin{frame}{Personal Area Network}
  \begin{block}{Problemas}    
     Além de alguns problemas já estarem presentes na BAN (consumo, segurança ...)
  \end{block}

  \begin{block}{A auto-organização dos nós é essencial}
    No sistema miTag ~\cite{ter08}, o dado coletado de uma rede sensor auto-organizável a qual opera em \text{mesh mode} transferindo dados para a internet. 

    O uso da operação em \textit{mesh mode} permite:
      \begin{itemize}
       \item extensibilidade
       \item cobertura
       \item distribuir processamento
       \item recuperação de desastres
       \item trabalhar em ambiente hostil
      \end{itemize}
  \end{block}
\end{frame}

% \section{Gateway para WAN}
\begin{frame}{Gateway para WAN}
  \begin{block}{}
    O gateway é um subsistema responsável em conectar a BAN ou PAN em WANs. Um gateway pode ser um aparelho de celular, PDA, nó sensor no ambiente assim como um laptop ou servidor.
  \end{block}

  \begin{block}{}
    A principal função de um gateway é prover uma conexão entre a rede sensor \textit{ad-hoc} e uma infraestrutura baseada numa WAN.
  \end{block}
\end{frame}

\begin{frame}{Gateway para WAN}
  \begin{block}{Problemas de Segurança}
   Os problemas de segurança também precisam ser tratados pelo \texit{gateway} do subsistema. Isso inclui: verficar a identidade da origem do dado, não motificar os dados do paciente, exceto se houver agregação de valor.
  \end{block}
\end{frame}

\section{Projetos}
\subsection{Monitoramento das atividades diárias}
\begin{frame}{Usuários Finais das Aplicações de Monitoramento}
  \begin{block}{}
   O ponto principal de todo esse arcabouço tecnológico é conseguir coletar os dados, interpretar e tomar as decisões corretas. Como por exemplo entender o batimento cardíaco, algoritimos de aprendizagem para identificar situações inesperadas, etc. 
  \end{block}
\end{frame}

% 
% \subsection{Segurança}
% \subsection{Privacidade}
% \subsection{Recuperação a Desastres}
% \subsection{Middleware Design}
% \subsection{Seamless healthcare tracking and monitoring system}

% \section{Aplicações de Monitoramento Healthcare e protótipos WSN}
% \subsection{}
\begin{frame}{Monitoramento \textit{Healthcare} e protótipos WSN}
  \begin{block}{}
    Existem algumas aplicações protótipos e comerciais para monitoramento de crianças, idosos e pacientes crônicos que incluem:
      \begin{itemize}[<+->]
       \item Monitoramento das atividades diárias;
       \item detecção de movimento e quedas;
       \item rastreamento;
       \item ingestão e monitoramento de medicamentos;
       \item status clínico.
      \end{itemize}
  \end{block}
\end{frame}


\begin{frame}{Monitoramento das atividades diárias}
  \begin{block}{Activities of Daily Living (ADL)}
    Lu e Fu ~\cite{lu} apresentam uma abordagem de reconhecimento de atividade baseada no contexto e localização. Eles usam uma rede sensor monitorando um ambiente inteligente integrado aos objetos domésticos.
  \end{block}
\end{frame}


\begin{frame}{Monitoramento das atividades diárias} 
    \begin{block}{Activities of Daily Living (ADL)}
      \center \includegraphics[height=2.4 in]{../imagens/hierarqlocat.png}
    \end{block}
\end{frame}

\begin{frame}{Monitoramento das atividades diárias}
  \begin{block}{Activities of Daily Living (ADL)}<1->
    Outra aplicação de monitoramento de atividade é o \textit{Caregiver’s Assistant} ~\cite{phi04}, desenvolvido para monitorar idosos em suas casas. Basicamente, o sistema funciona etiquetando com RFID diversos equipamentos (caixas de remédio, escovas de dentes, chaves, etc). 
  \end{block}

  \begin{block}{}<2->
   As etiquetas permitem identificar o que o usuário está utilizando naquele momento. Permitindo que os médicos possam identificar desvios de padrões das atividades normais, para verificar os níveis de demência ou desvio das atividades normais dos idosos.
  \end{block}
\end{frame}

\subsection{Detecção de movimentos e quedas}
\begin{frame}{Aplicações de detecção de movimentos e quedas}
  \begin{block}{}
     Enquanto que a detecção de movimentos e quedas são casos específicos de classificação das atividades, existe pesquisas significantes e esforços focados na detecção de quedas e postura corpórea e análise de locomoção.
  \end{block}

  \begin{alertblock}{}
     Acidentes provenientes de queda lideram as causas de morte de idosos (acima dos 65 anos de idade)
  \end{alertblock}
\end{frame}

\begin{frame}{Aplicações de detecção de movimentos e quedas}
  \begin{block}{}
    O projeto Information Technology for Assisted Living at Home
(ITALH) ~\cite{han05}, propõe habilitar o vídeo da câmera de celular somente em condições de emergência. Através de um sistema de detecção de quedas baseado no acelerômetro do celular em idosos. O nó sensor é capaz de analisar o acelerômetro em tempo real e classificar os eventos como variações normais(sentar, deitar, etc) ou anormais (quedas).
  \end{block}
\end{frame}

\begin{frame}{Hipguard} 
    \begin{block}{ADL de Análise de Postura}
      O HipGuard ~\cite{ket08} é um sistema de deteção de postura, destinado ao período de recuperação de 8-12 semanas após uma operação no quadril. O sistema tem sete nós sensores: na cintura, coxas, canelas e pés.     
    \end{block}

    \begin{block}{}
      A unidade de controle central coleta e processa os dados dos nós sensores e produz alarmes como sinal de áudio ou uma vibração háptica se a posição do quadril operado ou a carga colocada sobre os perna operada se aproxima dos limites estabelecidos pelos profissionais de saúde.
    \end{block}
\end{frame}

\begin{frame}{Hipguard} 
    \begin{block}{ADL de Análise de Postura}
      \center \includegraphics[height=2.4 in]{../imagens/hipguard.png}
    \end{block}
\end{frame}

\subsection{Monitoramento de ingestão de medicamentos}
\begin{frame}{Aplicações de Monitoramento de Ingestão de Medicamentos}
  \begin{block}{}<1->
   A não aderência da ingestão de medicamentos em idosos e pacientes crônicos é bastante freqüente. Portanto, o monitoramento e controle da ingestão da medicação torna-se essencial para um melhor acompanhamento do tratamento.
  \end{block}

  \begin{block}{}<2->
    Um protótipo inicial desenvolvido por Moh et al. ~\cite{Ho:2005:PRS:1080148.1080164} permite o controle de ingestão de medicamento para idosos, combinando redes sensores sem fio e RFID.
  \end{block}
\end{frame}

\begin{frame}{Protótipo de RFID com RSSF para Ingestão de Medicamentos em Idosos}
  \begin{block}{}
  Esta aplicação consiste de três subsistemas: identificação do vasilhame, etiqueta RFID identificando o que foi ingerido pelo paciente medido por uma escala de peso.
  \end{block}
\end{frame}

\begin{frame}{Projeto \textit{iCabinet}}
  \begin{block}{}
   O Projeto iCabinet ~\cite{lop08} faz uso do RFID inteligentes nas embalagem de comprimidos que podem gravar a remoção de um comprimido simplesmente quebrando um fluxo elétrico no circuito integrado do RFID. iCabiNET é uma aplicação interna que também está com interface para uma rede residencial, desta forma, a ingestão de medicamento pode ser monitorada através da rede residencial por intermédio dos leitores RFID presentes na casa.
  \end{block}
\end{frame}

\begin{frame}{Projeto \textit{iCabinet}}
  \begin{block}{Outras funções}
   \begin{itemize}[<+->]
    \item Monitoramento dos remédios adquiridos pelo paciente;
    \item integração com a TV, provendo alarmes para ingestão de medicamentos;
    \item integração com a TV para aquisição de novos equipamentos;
    \item integração com o celular para lembrar o paciente para tomar a medicação.
   \end{itemize}
  \end{block}
\end{frame}

\begin{frame}{Projeto \textit{iPackage}}
  \begin{block}{}
    O iPackage, é diferente de RFID inteligentes. Esse projeto usa uma matriz de filmes de laminados (CDM), identificando qual comprimido foi aberto.
  \end{block}

  \begin{block}{iPackage}
      \center \includegraphics[height=1 in]{../imagens/intake.png}
  \end{block}
\end{frame}

\subsection{Status clínico através de monitoramento}
\begin{frame}{Status clínico através de monitoramento}
  \begin{block}{}
      Monitorar o estado de saúde de pacientes é uma das aplicações mais estudadas em computação pervasiva aplicada à saúde. Os sinais vitais são comumente usados através de ECG, pulso, oxímetros, freqüencia cardíaca, pressão arterial e etc.   
  \end{block}

  \begin{alertblock}{}
    A maioria dos estudos focam em capturar e enviar o dado para outra avalição remota.
  \end{alertblock}
\end{frame}

\begin{frame}{Mobihealth}
  \begin{block}{}
  O projeto Mobihealth ~\cite{mobhealth} é um dos primeiros a integrar vários dispositivos sensores vestíveis (Wearable).
    Os sensores medem e transmitem dados fisiólogicos continuamente, juntamente com áudio e vídeo disponibilizando as gravações para os prestadores de serviços de saúde responsáveis por porporcionar uma assistência remota rápida e confiável em caso de acidente.
  \end{block}
\end{frame}

\begin{frame}{MobiHealth}
  \begin{block}{Convergência de Tecnologias}
    O MobiHealth é importante por ser um dos primeiros estudos propor a convergência dos sistemas de rede diferentes como :
    \begin{itemize}
     \item BAN
     \item PAN 
     \item WAN 
    \end{itemize}
  \end{block}
  \begin{block}{}
   Com esse convergência foi possível habilitar uma saúde personalizada e móvel.
  \end{block}
\end{frame}


\begin{frame}{Codeblue}
  \begin{block}{}
  O projeto Codeblue ~\cite{codeblue} é uma plataforma de hardware e software. O hardware consiste de um oxímetro, dois ECG e um sensor de movimento. O principal objetivo dessa plataforma é prover coordenação e comunicação entre os equipamentos wireless médicos de forma ad hoc. Provendo mecanismos de descobrimento de sensores e os canais de comunicação adequados.
  \end{block}
\end{frame}


\begin{frame}{Lifeguard}
  \begin{block}{}
  O projeto Lifeguard inicialmente foi desenvolvido para os astronautas, e pode ser usado para o monitoramento dos sinais vitais.
  A estação base parte é um Bluetooth Tablet PC que mostra e armazena os dados para avaliação futura.
  \end{block}
\end{frame}

\begin{frame}{Lifeguard}
  \begin{block}{Lifeguard}
      \center \includegraphics[height=2 in]{../imagens/lifeguard.png}
  \end{block}
\end{frame}


\section{Problemas em Aberto}
\subsection{}
\begin{frame}{Solução Não-Invasiva}
  \begin{block}{}
    \begin{itemize}[<+->]
     \item A concepção de um sistema de monitoramento vestível e que não seja invasivo ainda é um grande desafio; 
     \item a necessidade de integrar diferentes sensores em uma única solução torna essa atividade mais difícil;
     \item alguns desses dispositivos são pesados \textbf{visíveis e estereotipados}, como um ecg por exemplo ~\cite{aarh10}.
    \end{itemize}
  \end{block}
\end{frame}

\begin{frame}{Precisão e Calibragem}
  \begin{block}{}    
    A precisão dos sensores é de extrema importância, especialmente quando os usuários estão com seus sensores em ambientes hostis ou praticando exercícios. O suor pode afetar negativamente na precisão dos dados requerendo uma recalibração dos sensores.

    Gietzelt ~\cite{giet08}, propõe um algoritmo de calibração de acelerômetros triaxiais.
  \end{block}
\end{frame}

\begin{frame}{Energia}
  \begin{block}{}    
    O problema de eficiência energética é um dos principais problemas de RSSF. Para pacientes que têm acesso a energia pode-se utilizar baterias recarregáveis. Mas isso pode ser um problema comprometendo todo o monitoramento principalmente para idosos. Baixo consumo e um possível uso da energia solar pode ser uma das soluções.
  \end{block}
\end{frame}

\begin{frame}{Eficiência na Aquisição dos Dados}
  \begin{block}{}    
   É de extrema importância coletar os dados em sistemas de \textit{pervasive healthcare}. O desenvolvimento de técnicas de processamento desses dados é de suma importância. Em alguns casos 3 ECG não são suficientes para identificar uma doença cardíaca, ou 3 acelerômetros podem não ser capaz de identificar todas as atividades de uma pessoa. Nesses casos mais sensores são necessários e o conjunto de dados irá aumentar.
  \end{block}
\end{frame}

\begin{frame}{Eficiência na Aquisição dos Dados}
  \begin{block}{}    
  A aquisição e análise em tempo real dos dados fisiológicos é essencial. Além disso, o time-stamp e a ordenação dos eventos, sincronização dos diferentes sensores são problemas em aberto. Finalmente integrar diferentes tipos de sensores como: RFID, sensores implantados e sensores wirelesse necessitam do desenvolvimento de uma arquitetura modular.
  \end{block}
\end{frame}

\begin{frame}{Interoperabilidade}
  \begin{block}{}
    A integração de diversos tipos de sensores operando em freqüencias distintas elevam os problemas de interoperabilidade. A comunicação entre os diversos tipos de devices ocupando múltiplas bandas e usando os diferentes protocolos. Essa situação pode causar interferência, logo os sistemas de \textit{pervasive healthcare} devem ser interoperáveis entre os diferentes dispositivos.
  \end{block}
\end{frame}

\begin{frame}{Largura de Banda}
  \begin{block}{}
  A largura de banda disponível para comunicação numa Wireless Body Network é relativamente baixa. Devido a necessidade de baixo-consumo se os sensores se comunicarem em ciclos de 10$\%$ do seu tempo. Isso pode ser um complicador se estiver relacionado com o processamento de imagens o qual requer níveis de comunição na ordem do Mbit/s. Por esse motivo há a necessidade de algoritmos de compreessão para transmissão de dados multimídia. Levando em consideração que os nós sensores têm a capacidade de processamento bastante limitada, isso limita bastante o seu uso sendo um impreditivo para as tecnologias atuais.
  \end{block}
\end{frame}

\begin{frame}{Problemas no Protocolo de Acesso ao Meio (MAC)}
  \begin{block}{}
  Monitoramento de \textit{healthcare} requerem situações de emergência a eventos e reportagem de dados periodicamente. Sobre situações de emergência, o dado deve ser garantido e entregue dentro de um delay adequado. Para esse propósito, os dados em momentos de emergência devem ser priorizados.   
  \end{block}
\end{frame}


\begin{frame}{Segurança}
  \begin{block}{}<1->
  Os requisitos de segurança são fundamentais para qualquer sistema, confidencialidade, integridade de dados, disponibilidade e controle de acessos. Para assegurar esses requisitos, métodos de criptografia podem ser usados para solucionar esses problemas.
  \end{block}
  
  \begin{block}{}<2->
   No trabalho de Andersen ~\cite{and09}, é proposto um mecanismo de criptografia baseada em curvas elípticas para a distribuição de chaves, buscando reduzir o consumo de energia.
  \end{block}

  \begin{block}{}<3->
  Outro grande problema é quando o paciente está inconsciente, onde não é possível obter a senha nesses casos os dados biométricos podem ser utilizados.
  \end{block}
\end{frame}


\begin{frame}{Privacidade}
  \begin{block}{}<1->
    Autorização do usuário no sistema não pode ser sobreposta dando total autonomia para o usuário controlar seus próprios dados.
  \end{block}
  
  \begin{block}{}<2->
   Aplicações de processamento de imagens devem prover o anônimato para os usuários.
  \end{block}

  \begin{block}{}<3->
  Sistemas de ADL devem ter monitores anti-ataques às redes wireless.
  \end{block}
\end{frame}

\begin{frame}{Facilidade ao Usuário}
 \begin{block}{}
  Sabendo que o sistema será utilizado, por idosos, pessoas com deficiência cognitiva, médicos, enfermeiros e pessoas com os diferentes perfis. É bastante complexo conseguir uma aplicação \textit{user friendly}.
 \end{block}
\end{frame}

\begin{frame}{Fácil Implantação e Escalabilidade}
  \begin{block}{}
  Similar a facilidade ao usuário, desenvolver uma aplicação de fácil implantação para sistemas pervasivos é tão essencial como não-trivial. As plataformas de software distribuídas precisam monitorar e integrar hardware e software.
  \end{block}
\end{frame}

\section{Conclusão}
\subsection{Conclusão}
\begin{frame}{Conclusão}
  \begin{block}{}
    O principal objetivo de monitorar a saúde é permitir que as pessoas sobrevivam de forma independente e com uma boa qualidade de vida através dos serviço de healthcare. O uso de sensores e redes-sensores com esse propósito não é novo, entretanto com o surgimento das redes sensor sem-fio foi possível o desenvolvimento de aplicações móveis para os usuários.
  \end{block}

  \begin{block}{}
   As redes sensores podem prover a mobilidade e o monitoramento constante dos pacientes de forma ubíqua e pervasiva, possibilitando inclusive o desenvolvimento de aplicação baseadas em localidade e sensíveis ao contexto.
  \end{block}
\end{frame}


\subsection{Dúvidas}
\begin{frame}
  \begin{center}
  DÚVIDAS ?
  \end{center}
\end{frame}

\subsection{Bibliografia}
\bibliographystyle{alpha}
%\bibliographystyle{plain} % estilo de bibliografia   plain,unsrt,alpha,abbrv.
%\bibliography{../../tex/bibtex/biblmm} % arquivos com as entradas bib.
\bibliography{wsnbibtex} % arquivos com as entradas bib.

\end{document}
